% !TEX spellcheck = en_US



% ~~~~~~~~~~~~~~~~~~~~~~~~~~~~~~~~~~~~~~~ V
% (NC) 2026 Haluk Bingol
% github.com/halukbingol/LaTeX-Templates
%
% Licensees may copy, distribute, display, and perform the work and make derivative works 
% and remixes based on it only for non-commercial purposes. 
% ~~~~~~~~~~~~~~~~~~~~~~~~~~~~~~~~~~~~~~~ A




% =====================================================================
%  Trees -- Senior Level: Modern Java, Persistence, Parallelism and Real Trees
%  ---------------------------------------------------------------------
%  Built on hblib/Template/hblibTemplate.tex with the libraries in ../../hblib/.
%  Programs and their real outputs are in codeoutput/:
%     codeoutput/<Name>.java  a program (classes without main() are only compiled)
%     codeoutput/<Name>.sh    a terminal session, run in a copy of codeoutput/<Name>/
%     codeoutput/<Name>.txt   what it printed, made by:  bash run_examples.sh
% =====================================================================

\documentclass[aspectratio=169,10pt,compress]{beamer}

% ~~~~~~~~~~~~~~~~~~~~~~~~~~~~~~~~~~~~~~~~~~~~~~~~~~~~~~~~~~~~~~~~~ V hblib
	\usepackage{../../hblib/hblibPresentation} % lib presentation: theme, i/N, \hSection, algorithm2e, ...
	\usepackage{../../hblib/hblibCodeoutput} % lib code + output: \lstcode, \lstterminal, ...
	\usepackage{../../hblib/hblibDatastructures} % lib data structures: \hString, \hStack, \hQueue
% ~~~~~~~~~~~~~~~~~~~~~~~~~~~~~~~~~~~~~~~~~~~~~~~~~~~~~~~~~~~~~~~~~ A hblib


% xcolor, array (and the L{width} column), listings, tikz and algorithm2e come from the hblib libraries
\usepackage[T1]{fontenc}
\usepackage{lmodern}
\usepackage{booktabs}

\newcommand{\win}{\textbf{Windows}}
\newcommand{\mac}{\textbf{macOS}}

% algorithm2e set-up (the package is loaded by hblibPresentation)
\SetAlCapFnt{\small}
\SetAlCapNameFnt{\small}
\SetKwProg{Fn}{function}{}{end}
\SetKw{KwError}{error}

% ---------------------------------------------------------------------
\title{Trees}
\subtitle{Senior Level: Modern Java, Persistence, Parallelism and Real Trees}
\author[Bingol]{Haluk O. Bingol}

\institute[FCIS]{
    Faculty of Computer and Information Sciences\\
    Yeditepe University
}
\date{
	{\tiny \hbTimeStamp}
}

\begin{document}




% ~~~~~~~~~~~~~~~~~~~~~~~~~~~~~~~~~~~~~~ V frm
\begin{frame}[t,fragile]
  \titlepage
\end{frame}
% ~~~~~~~~~~~~~~~~~~~~~~~~~~~~~~~~~~~~~~ A frm




% ~~~~~~~~~~~~~~~~~~~~~~~~~~~~~~~~~~~~~~ V frm
\begin{frame}[t,fragile] \frametitle{Outline}
  \begin{columns}[T]
    \begin{column}{0.48\textwidth}
      \tableofcontents[sections={1-3}, hideallsubsections]
    \end{column}
    \begin{column}{0.48\textwidth}
      \tableofcontents[sections={4-6}, hideallsubsections]
    \end{column}
  \end{columns}
\end{frame}
% ~~~~~~~~~~~~~~~~~~~~~~~~~~~~~~~~~~~~~~ A frm

% ~~~~~~~~~~~~~~~~~~~~~~~~~~~~~~~~~~~~~~ V
\hSection[Modern Java]{Trees in Modern Java}
% ~~~~~~~~~~~~~~~~~~~~~~~~~~~~~~~~~~~~~~ A




% ~~~~~~~~~~~~~~~~~~~~~~~~~~~~~~~~~~~~~~ V frm
\begin{frame}[t,fragile] \frametitle{Sealed types and records: an algebraic data type}
  \begin{columns}[T]
    \begin{column}{0.49\textwidth}
      \lstcode[basicstyle=\ttfamily\tiny, lastline=7]{codeoutput/Expr.java}{Java}
    \end{column}
    \begin{column}{0.49\textwidth}
      \lstcode[basicstyle=\ttfamily\tiny, firstline=8, firstnumber=8, lastline=25]{codeoutput/Expr.java}{Java}
    \end{column}
  \end{columns}
\end{frame}
% ~~~~~~~~~~~~~~~~~~~~~~~~~~~~~~~~~~~~~~ A frm




% ~~~~~~~~~~~~~~~~~~~~~~~~~~~~~~~~~~~~~~ V frm
\begin{frame}[t,fragile] \frametitle{What the new Java features do}
  \begin{columns}[T]
    \begin{column}{0.36\textwidth}
      \begin{center}
      \begin{tikzpicture}[htree, level distance=0.95cm, level 1/.style={sibling distance=2.8cm}, level 2/.style={sibling distance=1.5cm}, level 3/.style={sibling distance=0.8cm}, scale=0.7, transform shape]
  \node[hsquare] {Add} child {node[hsquare] {Mul} child {node[hnode] {3}} child {node[hsquare] {Mul} child {node[hnode] {x}} child {node[hnode] {x}}}} child {node[hsquare] {Add} child {node[hsquare] {Mul} child {node[hnode] {2}} child {node[hnode] {x}}} child {node[hnode] {1}}};
\end{tikzpicture}\\
      \small $3x^{2} + 2x + 1$ as \texttt{Expr} records
      \end{center}
    \end{column}
    \begin{column}{0.6\textwidth}
      \begin{itemize}\small
        \item \texttt{sealed \dots{} permits}: these four records are \textbf{all} the node kinds (Java 17).
        \item Records: immutable nodes with \texttt{equals}, \texttt{hashCode} and \texttt{toString} for free (Java 16).
        \item \texttt{case Add(Expr l, Expr r)}: a \textbf{record pattern} takes a node apart (Java 21).
        \item No \texttt{default}: a \texttt{switch} that forgets a node kind does not \textbf{compile}, so adding a node kind shows every place to update.
      \end{itemize}
    \end{column}
  \end{columns}
\end{frame}
% ~~~~~~~~~~~~~~~~~~~~~~~~~~~~~~~~~~~~~~ A frm




% ~~~~~~~~~~~~~~~~~~~~~~~~~~~~~~~~~~~~~~ V frm
\begin{frame}[t,fragile] \frametitle{Transforming trees: derivative and simplification}
  \begin{columns}[T]
    \begin{column}{0.49\textwidth}
      \lstcode[basicstyle=\ttfamily\tiny, firstline=26, firstnumber=26, lastline=34]{codeoutput/Expr.java}{Java}
    \end{column}
    \begin{column}{0.49\textwidth}
      \lstcode[basicstyle=\ttfamily\tiny, firstline=35, firstnumber=35]{codeoutput/Expr.java}{Java}
    \end{column}
  \end{columns}
\end{frame}
% ~~~~~~~~~~~~~~~~~~~~~~~~~~~~~~~~~~~~~~ A frm




% ~~~~~~~~~~~~~~~~~~~~~~~~~~~~~~~~~~~~~~ V frm
\begin{frame}[t,fragile] \frametitle{Symbolic differentiation}
  \lstcoderightstep[0.6][basicstyle=\ttfamily\tiny][basicstyle=\ttfamily\tiny]{codeoutput/UseExpr.java}{Java}{codeoutput/UseExpr.txt}

  \vspace{0.2em}
  \small \texttt{derive} builds a new tree in one traversal; \texttt{simplify} works \textbf{postorder}: children are simplified before their parent is examined.
\end{frame}
% ~~~~~~~~~~~~~~~~~~~~~~~~~~~~~~~~~~~~~~ A frm




% ~~~~~~~~~~~~~~~~~~~~~~~~~~~~~~~~~~~~~~ V frm
\begin{frame}[t,fragile] \frametitle{Pattern matching or the Visitor pattern?}
  \begin{columns}[T]
    \begin{column}{0.54\textwidth}


% +++++++++++++++++++++++++++++++++++++++ V lst
%: -Lst.
\begin{lstlisting}[style=codesnippet, language=Java, basicstyle=\ttfamily\scriptsize]
interface Visitor<R> {
    R visitNum(Num n);
    R visitAdd(Add a);
    R visitMul(Mul m);
}
interface Expr { <R> R accept(Visitor<R> v); }
record Add(Expr l, Expr r) implements Expr {
    public <R> R accept(Visitor<R> v) { return v.visitAdd(this); }
}
\end{lstlisting}
% +++++++++++++++++++++++++++++++++++++++ A lst


    \end{column}
    \begin{column}{0.42\textwidth}
      \begin{itemize}\small
        \item The classic \textbf{Visitor} (Gang of Four, 1994) adds operations without changing the node classes, through double dispatch.
        \item With sealed types and pattern matching, each operation is \textbf{one method with a switch}: shorter, checked for completeness.
        \item Visitors are still common in older code bases: \texttt{javax.lang.model}, ASM, ANTLR, \texttt{Files.walkFileTree}.
      \end{itemize}
    \end{column}
  \end{columns}
\end{frame}
% ~~~~~~~~~~~~~~~~~~~~~~~~~~~~~~~~~~~~~~ A frm

% ~~~~~~~~~~~~~~~~~~~~~~~~~~~~~~~~~~~~~~ V
\hSection[Persistence]{Immutable and Persistent Trees}
% ~~~~~~~~~~~~~~~~~~~~~~~~~~~~~~~~~~~~~~ A




% ~~~~~~~~~~~~~~~~~~~~~~~~~~~~~~~~~~~~~~ V frm
\begin{frame}[t,fragile] \frametitle{Path copying}
  \begin{columns}[T]
    \begin{column}{0.49\textwidth}
      \lstcode[basicstyle=\ttfamily\tiny, lastline=21]{codeoutput/Persistent.java}{Java}
    \end{column}
    \begin{column}{0.49\textwidth}
      \lstcode[basicstyle=\ttfamily\tiny, firstline=22, firstnumber=22]{codeoutput/Persistent.java}{Java}
    \end{column}
  \end{columns}
\end{frame}
% ~~~~~~~~~~~~~~~~~~~~~~~~~~~~~~~~~~~~~~ A frm




% ~~~~~~~~~~~~~~~~~~~~~~~~~~~~~~~~~~~~~~ V frm
\begin{frame}[t,fragile] \frametitle{Path copying: the result}
  \begin{columns}[T]
    \begin{column}{0.5\textwidth}
      \begin{center}
      \begin{tikzpicture}[htree, level distance=0.95cm, level 1/.style={sibling distance=2.6cm}, level 2/.style={sibling distance=1.3cm}, level 3/.style={sibling distance=0.65cm}, scale=0.8, transform shape]
  \node[hnode, hhot] {50} child {node[hnode] {30} child {node[hnode] {20} child {node[hnode] {10}} child {node[hnode] {25}}} child {node[hnode] {40} child {node[hnode] {35}} child {node[hnode] {45}}}} child {node[hnode, hhot] {70} child {node[hnode, hhot] {60} child[missing] child {node[hnode, hhot] {65}}} child {node[hnode] {80}}};
\end{tikzpicture}
      \end{center}
      \small green: new nodes in \texttt{v2}; white: shared with \texttt{v1}
    \end{column}
    \begin{column}{0.46\textwidth}
      \lstoutput[basicstyle=\ttfamily\tiny]{codeoutput/Persistent.txt}
      \begin{itemize}\small
        \item Only the \textbf{path} from the root to the new key is copied: $O(h)$ new nodes; the rest is shared.
        \item \texttt{v1} is unchanged and still usable: every version is \emph{persistent}.
        \item Safe to share between threads without locks; undo for free; used in functional languages, Git, and databases (copy-on-write B-trees).
      \end{itemize}
    \end{column}
  \end{columns}
\end{frame}
% ~~~~~~~~~~~~~~~~~~~~~~~~~~~~~~~~~~~~~~ A frm

% ~~~~~~~~~~~~~~~~~~~~~~~~~~~~~~~~~~~~~~ V
\hSection[Parallel]{Parallel Traversal}
% ~~~~~~~~~~~~~~~~~~~~~~~~~~~~~~~~~~~~~~ A




% ~~~~~~~~~~~~~~~~~~~~~~~~~~~~~~~~~~~~~~ V frm
\begin{frame}[t,fragile] \frametitle{Fork/join: sum a tree in parallel}
  \begin{columns}[T]
    \begin{column}{0.49\textwidth}
      \lstcode[basicstyle=\ttfamily\tiny, lastline=15]{codeoutput/ParallelSum.java}{Java}
    \end{column}
    \begin{column}{0.49\textwidth}
      \lstcode[basicstyle=\ttfamily\tiny, firstline=16, firstnumber=16, lastline=29]{codeoutput/ParallelSum.java}{Java}
    \end{column}
  \end{columns}
\end{frame}
% ~~~~~~~~~~~~~~~~~~~~~~~~~~~~~~~~~~~~~~ A frm




% ~~~~~~~~~~~~~~~~~~~~~~~~~~~~~~~~~~~~~~ V frm
\begin{frame}[t,fragile] \frametitle{Fork/join: timing both versions}
  \lstcode[basicstyle=\ttfamily\scriptsize, firstline=30, firstnumber=30]{codeoutput/ParallelSum.java}{Java}
\end{frame}
% ~~~~~~~~~~~~~~~~~~~~~~~~~~~~~~~~~~~~~~ A frm




% ~~~~~~~~~~~~~~~~~~~~~~~~~~~~~~~~~~~~~~ V frm
\begin{frame}[t,fragile] \frametitle{Fork/join: the result}
  \begin{columns}[T]
    \begin{column}{0.5\textwidth}
      \lstoutput[basicstyle=\ttfamily\scriptsize]{codeoutput/ParallelSum.txt}
    \end{column}
    \begin{column}{0.46\textwidth}
      \begin{itemize}\small
        \item The two subtrees are independent: \texttt{fork()} one, compute the other, \texttt{join()}.
        \item Below depth 6 the work is done sequentially: tasks must not be too small.
        \item Times are from one run. On this 2-processor machine the parallel sum was \textbf{not} faster: adding numbers is too little work per node, so memory, not the processors, is the limit.
        \item \texttt{availableProcessors()} counts the processors the JVM may use, e.g.\ in a container.
      \end{itemize}
    \end{column}
  \end{columns}
\end{frame}
% ~~~~~~~~~~~~~~~~~~~~~~~~~~~~~~~~~~~~~~ A frm

% ~~~~~~~~~~~~~~~~~~~~~~~~~~~~~~~~~~~~~~ V
\hSection[Library]{Trees in the Java Library}
% ~~~~~~~~~~~~~~~~~~~~~~~~~~~~~~~~~~~~~~ A




% ~~~~~~~~~~~~~~~~~~~~~~~~~~~~~~~~~~~~~~ V frm
\begin{frame}[t,fragile] \frametitle{\texttt{TreeMap} as a \texttt{NavigableMap}}
  \lstcoderightstep[0.62][basicstyle=\ttfamily\tiny][basicstyle=\ttfamily\tiny]{codeoutput/Navigable.java}{Java}{codeoutput/Navigable.txt}

  \vspace{0.2em}
  \small A red-black tree answers ``nearest key'' questions in $O(\log n)$: \texttt{floorEntry} turns a score into a letter grade without any \texttt{if}.
\end{frame}
% ~~~~~~~~~~~~~~~~~~~~~~~~~~~~~~~~~~~~~~ A frm




% ~~~~~~~~~~~~~~~~~~~~~~~~~~~~~~~~~~~~~~ V frm
\begin{frame}[t,fragile] \frametitle{Which sorted structure?}
  \begin{center}\small
  \begin{tabular}{L{0.38\textwidth}L{0.3\textwidth}L{0.22\textwidth}}
    \toprule
    \textbf{class}                    & \textbf{structure}                    & \textbf{use} \\
    \midrule
    \texttt{TreeMap}, \texttt{TreeSet} & red-black tree                       & sorted, single thread \\
    \texttt{ConcurrentSkipListMap}    & skip list (not a tree)                & sorted, many threads, no locks \\
    \texttt{Collections.synchronizedSortedMap} & a \texttt{TreeMap} with one lock & simple, but every call locks \\
    \texttt{PriorityQueue}            & binary heap in an array               & only the smallest element \\
    \texttt{HashMap}                  & hash table; a bin with 8 or more keys becomes a red-black tree & no order, $O(1)$ \\
    \bottomrule
  \end{tabular}
  \end{center}
  \small Since Java 8, \texttt{HashMap} turns long collision chains into small trees, so many keys with the same hash cost $O(\log n)$ instead of $O(n)$: a defence against hash-flooding attacks.
\end{frame}
% ~~~~~~~~~~~~~~~~~~~~~~~~~~~~~~~~~~~~~~ A frm

% ~~~~~~~~~~~~~~~~~~~~~~~~~~~~~~~~~~~~~~ V
\hSection[Real trees]{Real Trees: Documents and Databases}
% ~~~~~~~~~~~~~~~~~~~~~~~~~~~~~~~~~~~~~~ A




% ~~~~~~~~~~~~~~~~~~~~~~~~~~~~~~~~~~~~~~ V frm
\begin{frame}[t,fragile] \frametitle{XML is an N-ary tree: the DOM}
  \begin{columns}[T]
    \begin{column}{0.49\textwidth}
      \lstcode[basicstyle=\ttfamily\tiny, lastline=21]{codeoutput/Dom.java}{Java}
    \end{column}
    \begin{column}{0.49\textwidth}
      \lstcode[basicstyle=\ttfamily\tiny, firstline=22, firstnumber=22]{codeoutput/Dom.java}{Java}
    \end{column}
  \end{columns}
\end{frame}
% ~~~~~~~~~~~~~~~~~~~~~~~~~~~~~~~~~~~~~~ A frm




% ~~~~~~~~~~~~~~~~~~~~~~~~~~~~~~~~~~~~~~ V frm
\begin{frame}[t,fragile] \frametitle{The DOM: the result}
  \begin{columns}[T]
    \begin{column}{0.5\textwidth}
      \lstoutput[basicstyle=\ttfamily\scriptsize]{codeoutput/Dom.txt}
    \end{column}
    \begin{column}{0.46\textwidth}
      \begin{itemize}\small
        \item \texttt{getFirstChild} / \texttt{getNextSibling}: the DOM is stored as \textbf{first child, next sibling}.
        \item Text and whitespace are nodes too; the walk prints only elements. The \texttt{[Fatal Error]} line is printed by the parser itself.
        \item \hRemark{Security}: a \texttt{DOCTYPE} can make the parser read local files or expand entities exponentially (XXE, ``billion laughs). Disallow it for untrusted XML (OWASP).
      \end{itemize}
    \end{column}
  \end{columns}
\end{frame}
% ~~~~~~~~~~~~~~~~~~~~~~~~~~~~~~~~~~~~~~ A frm




% ~~~~~~~~~~~~~~~~~~~~~~~~~~~~~~~~~~~~~~ V frm
\begin{frame}[t,fragile] \frametitle{Why databases use trees with many children}
  A \textbf{B-tree} node holds hundreds of keys and children, sized to one disk or memory page. With fan-out $m$ and $n$ keys,
  \[
  	h \approx \log_{m} n = \frac{\log_{2} n}{\log_{2} m}
  \]
  \begin{center}\small
  \begin{tabular}{lccc}
    \toprule
    $n = 10^{9}$ keys          & binary ($m = 2$) & $m = 100$ & $m = 1000$ \\
    \midrule
    levels (page reads)        & 30               & 5         & 3 \\
    \bottomrule
  \end{tabular}
  \end{center}
  \begin{itemize}\small
    \item Reading a page costs far more than comparing keys inside it: minimise the \textbf{height}, not the comparisons.
    \item B+-trees (keys in the leaves, leaves linked for range scans) index almost every relational database:
          PostgreSQL, MySQL InnoDB, SQLite; file systems such as NTFS (\win), APFS (\mac), btrfs.
  \end{itemize}
\end{frame}
% ~~~~~~~~~~~~~~~~~~~~~~~~~~~~~~~~~~~~~~ A frm

% ~~~~~~~~~~~~~~~~~~~~~~~~~~~~~~~~~~~~~~ V
\hSection[Summary]{Summary}
% ~~~~~~~~~~~~~~~~~~~~~~~~~~~~~~~~~~~~~~ A




% ~~~~~~~~~~~~~~~~~~~~~~~~~~~~~~~~~~~~~~ V frm
\begin{frame}[t,fragile] \frametitle{Big ideas}
  \begin{itemize}
    \item Sealed interfaces, records and pattern matching make tree code short and \textbf{checked} by the compiler.
    \item Immutable trees with path copying give persistent versions in $O(h)$ extra nodes each.
    \item Independent subtrees can be processed in parallel with fork/join; keep tasks large enough.
    \item \texttt{TreeMap} is a red-black tree with a rich navigation API; \texttt{HashMap} uses trees inside bins.
    \item Real data is tree-shaped (XML, JSON, file systems); parse untrusted documents defensively.
    \item High fan-out trees (B-trees) minimise page reads: databases and file systems.
  \end{itemize}
\end{frame}
% ~~~~~~~~~~~~~~~~~~~~~~~~~~~~~~~~~~~~~~ A frm




% ~~~~~~~~~~~~~~~~~~~~~~~~~~~~~~~~~~~~~~ V frm
\begin{frame}[t,fragile] \frametitle{Projects}
  \begin{enumerate}\small
    \item Extend \texttt{Expr} with \texttt{Sub}, \texttt{Div}, \texttt{Pow} and a parser from text; make \texttt{simplify} reach a fixed point.
    \item Make \texttt{Persistent} balanced (a persistent red-black or AVL tree) and keep 1000 versions; measure the shared nodes.
    \item Compare fork/join and sequential traversal for trees of $2^{10} \dots 2^{24}$ nodes; find the best cut-off depth.
    \item Implement a B-tree of order $m$ in memory with insert and search; measure height for $m = 3, 16, 128$.
    \item Walk a JSON document (e.g.\ with the \texttt{jakarta.json} API) and report its maximum depth; reject documents deeper than 100.
  \end{enumerate}
\end{frame}
% ~~~~~~~~~~~~~~~~~~~~~~~~~~~~~~~~~~~~~~ A frm




% ~~~~~~~~~~~~~~~~~~~~~~~~~~~~~~~~~~~~~~ V frm
\begin{frame}[t,fragile]
	\begin{center}
		\vspace{4cm}
		\hRemark{\Huge Questions?}
	\end{center}
\end{frame}
% ~~~~~~~~~~~~~~~~~~~~~~~~~~~~~~~~~~~~~~ A frm



\end{document}
